\documentclass[12pt,a4paper]{article}

% ---- Encoding & Font ----
\usepackage{iftex}
\ifPDFTeX
  \usepackage[utf8]{inputenc}
  \usepackage[T1]{fontenc}
\else
  \usepackage{fontspec}  % XeTeX/LuaTeX (tectonic): Unicode nativo (·, —, ¿, ª, §)
\fi
\usepackage[spanish]{babel}
\usepackage{csquotes}

% ---- Page layout ----
\usepackage[top=2.5cm, bottom=2.5cm, left=2.5cm, right=2.5cm]{geometry}
\usepackage{setspace}
\onehalfspacing

% ---- Math ----
\usepackage{amsmath, amssymb, amsthm}
\usepackage{mathtools}
\usepackage{physics}
\usepackage{esint}  % \oiint

% ---- Graphics ----
\usepackage{graphicx}
\usepackage{tikz}
\usepackage{float}
\usepackage{caption}

% ---- Tables ----
\usepackage{array}
\usepackage{booktabs}
\usepackage{tabularx}
\usepackage{colortbl}

% ---- Colours ----
\usepackage{xcolor}
\definecolor{notebg}{HTML}{FFF3CD}
\definecolor{noteborder}{HTML}{FFC107}
\definecolor{boxred}{HTML}{F8D7DA}
\definecolor{boxredborder}{HTML}{F5C6CB}

% ---- Boxes ----
\usepackage{mdframed}
\usepackage{tcolorbox}
\tcbuselibrary{most}

\newtcolorbox{keybox}[1][]{
  colback=blue!5,
  colframe=blue!70!black,
  arc=4pt,
  boxrule=1pt,
  fonttitle=\bfseries,
  title={#1}
}

\newtcolorbox{warnbox}[1][]{
  colback=boxred,
  colframe=boxredborder,
  coltitle=black!75,
  arc=4pt,
  boxrule=1pt,
  fonttitle=\bfseries,
  title={#1}
}

\newtcolorbox{notebox}[1][]{
  colback=notebg,
  colframe=noteborder,
  coltitle=black!75,
  arc=4pt,
  boxrule=1pt,
  fonttitle=\bfseries,
  title={#1}
}

% ---- Hyperlinks & TOC ----
\usepackage[hidelinks]{hyperref}
\usepackage{bookmark}

% ---- Enumerate ----
\usepackage{enumitem}

% ---- Miscellaneous ----
\usepackage{icomma}
\usepackage{siunitx}
\usepackage{textcomp}
\usepackage[bottom]{footmisc}

% ---- Diagramas (gráficas y esquemas) ----
\usepackage{pgfplots}
\pgfplotsset{compat=1.18}
\usetikzlibrary{babel}  % babel español activa `>`; esto evita romper las flechas `->`
% courses/ElecMag2024/Clases/assets/tikzstyles.tex
% ---------------------------------------------------------------------------
% Estilos compartidos para los diagramas del curso Electromagnetismo (2024).
% Fuente única del "look" visual (colores, grosores, escalas) para que todas
% las figuras (TikZ / pgfplots / CircuiTikz) sean consistentes entre clases.
%
% Es una COPIA de courses/Fisica3-2015/Clases/assets/tikzstyles.tex: mismo
% dominio, misma paleta. Copia y no symlink para que cada curso sea
% autocontenido y la paleta de Física III pueda evolucionar aparte.
%
% Uso: la clase debe cargar antes `tikz` y `pgfplots` (y `circuitikz` si la
%      clase tiene circuitos), y luego % courses/ElecMag2024/Clases/assets/tikzstyles.tex
% ---------------------------------------------------------------------------
% Estilos compartidos para los diagramas del curso Electromagnetismo (2024).
% Fuente única del "look" visual (colores, grosores, escalas) para que todas
% las figuras (TikZ / pgfplots / CircuiTikz) sean consistentes entre clases.
%
% Es una COPIA de courses/Fisica3-2015/Clases/assets/tikzstyles.tex: mismo
% dominio, misma paleta. Copia y no symlink para que cada curso sea
% autocontenido y la paleta de Física III pueda evolucionar aparte.
%
% Uso: la clase debe cargar antes `tikz` y `pgfplots` (y `circuitikz` si la
%      clase tiene circuitos), y luego % courses/ElecMag2024/Clases/assets/tikzstyles.tex
% ---------------------------------------------------------------------------
% Estilos compartidos para los diagramas del curso Electromagnetismo (2024).
% Fuente única del "look" visual (colores, grosores, escalas) para que todas
% las figuras (TikZ / pgfplots / CircuiTikz) sean consistentes entre clases.
%
% Es una COPIA de courses/Fisica3-2015/Clases/assets/tikzstyles.tex: mismo
% dominio, misma paleta. Copia y no symlink para que cada curso sea
% autocontenido y la paleta de Física III pueda evolucionar aparte.
%
% Uso: la clase debe cargar antes `tikz` y `pgfplots` (y `circuitikz` si la
%      clase tiene circuitos), y luego \input{../assets/tikzstyles.tex}
% (compilar desde el directorio de la clase para que la ruta relativa resuelva).
% ---------------------------------------------------------------------------

% ---- Paleta (alineada con las cajas del preámbulo: keybox/notebox/warnbox) ----
\definecolor{figblue}{HTML}{1F4E79}   % azul principal (líneas, keybox)
\definecolor{figred}{HTML}{C0392B}    % rojo de énfasis (warnbox)
\definecolor{figamber}{HTML}{B8860B}  % ámbar (notebox)
\definecolor{figgray}{HTML}{555555}   % auxiliares, ejes, guías

% ---- CircuiTikz: escala y grosor de los componentes ----
% Guarda \ifdefined: la electrostática (Clases 1-14) no carga `circuitikz`, y
% sin la guarda el \input revienta con `Undefined control sequence \ctikzset`.
% Cuando el curso llegue a circuitos, el bloque se activa solo.
\ifdefined\ctikzset
  \ctikzset{
    bipoles/thickness=1,
    resistors/scale=0.9,
    capacitors/scale=0.9,
  }
\fi

% ---- TikZ: estilos reutilizables ----
% Nota: las puntas se declaran como `-{Latex[...]}` y no como `->`. La forma
% con llaves NO contiene un `>` literal, así que es inmune al `>` activo que
% introduce babel español (ver CLAUDE.md §7.3.3).
\usetikzlibrary{arrows.meta,calc,angles,quotes}

\tikzset{
  figline/.style={figblue, line width=0.9pt},
  figaux/.style={figgray, dashed, line width=0.6pt},
  % vectores
  figvec/.style={figblue, line width=0.9pt, -{Latex[length=2.2mm,width=1.6mm]}},
  figvecres/.style={figred, line width=1.3pt, -{Latex[length=2.6mm,width=1.9mm]}},
  figvecaux/.style={figgray, line width=0.7pt, -{Latex[length=1.8mm,width=1.3mm]}},
  figaxis/.style={figgray, line width=0.6pt, -{Latex[length=2mm,width=1.4mm]}},
  % cargas puntuales
  figcarga/.style={circle, inner sep=0pt, minimum size=5.4pt, draw=none},
  figpos/.style={figcarga, fill=figred},
  figneg/.style={figcarga, fill=figblue},
  figneutra/.style={figcarga, fill=figgray},
  % etiquetas
  figlbl/.style={font=\small},
  figlblaux/.style={font=\footnotesize, figgray},
}

% ---- pgfplots: estilo base de las gráficas ----
\pgfplotsset{
  emfig/.style={
    width=11cm, height=6.5cm,
    axis lines=left,
    xlabel style={font=\small},
    ylabel style={font=\small},
    tick label style={font=\footnotesize},
    every axis plot/.append style={line width=1pt},
    grid=major,
    grid style={figgray!20},
    legend cell align=left,
    legend style={font=\footnotesize, draw=figgray!40},
  },
}

% (compilar desde el directorio de la clase para que la ruta relativa resuelva).
% ---------------------------------------------------------------------------

% ---- Paleta (alineada con las cajas del preámbulo: keybox/notebox/warnbox) ----
\definecolor{figblue}{HTML}{1F4E79}   % azul principal (líneas, keybox)
\definecolor{figred}{HTML}{C0392B}    % rojo de énfasis (warnbox)
\definecolor{figamber}{HTML}{B8860B}  % ámbar (notebox)
\definecolor{figgray}{HTML}{555555}   % auxiliares, ejes, guías

% ---- CircuiTikz: escala y grosor de los componentes ----
% Guarda \ifdefined: la electrostática (Clases 1-14) no carga `circuitikz`, y
% sin la guarda el \input revienta con `Undefined control sequence \ctikzset`.
% Cuando el curso llegue a circuitos, el bloque se activa solo.
\ifdefined\ctikzset
  \ctikzset{
    bipoles/thickness=1,
    resistors/scale=0.9,
    capacitors/scale=0.9,
  }
\fi

% ---- TikZ: estilos reutilizables ----
% Nota: las puntas se declaran como `-{Latex[...]}` y no como `->`. La forma
% con llaves NO contiene un `>` literal, así que es inmune al `>` activo que
% introduce babel español (ver CLAUDE.md §7.3.3).
\usetikzlibrary{arrows.meta,calc,angles,quotes}

\tikzset{
  figline/.style={figblue, line width=0.9pt},
  figaux/.style={figgray, dashed, line width=0.6pt},
  % vectores
  figvec/.style={figblue, line width=0.9pt, -{Latex[length=2.2mm,width=1.6mm]}},
  figvecres/.style={figred, line width=1.3pt, -{Latex[length=2.6mm,width=1.9mm]}},
  figvecaux/.style={figgray, line width=0.7pt, -{Latex[length=1.8mm,width=1.3mm]}},
  figaxis/.style={figgray, line width=0.6pt, -{Latex[length=2mm,width=1.4mm]}},
  % cargas puntuales
  figcarga/.style={circle, inner sep=0pt, minimum size=5.4pt, draw=none},
  figpos/.style={figcarga, fill=figred},
  figneg/.style={figcarga, fill=figblue},
  figneutra/.style={figcarga, fill=figgray},
  % etiquetas
  figlbl/.style={font=\small},
  figlblaux/.style={font=\footnotesize, figgray},
}

% ---- pgfplots: estilo base de las gráficas ----
\pgfplotsset{
  emfig/.style={
    width=11cm, height=6.5cm,
    axis lines=left,
    xlabel style={font=\small},
    ylabel style={font=\small},
    tick label style={font=\footnotesize},
    every axis plot/.append style={line width=1pt},
    grid=major,
    grid style={figgray!20},
    legend cell align=left,
    legend style={font=\footnotesize, draw=figgray!40},
  },
}

% (compilar desde el directorio de la clase para que la ruta relativa resuelva).
% ---------------------------------------------------------------------------

% ---- Paleta (alineada con las cajas del preámbulo: keybox/notebox/warnbox) ----
\definecolor{figblue}{HTML}{1F4E79}   % azul principal (líneas, keybox)
\definecolor{figred}{HTML}{C0392B}    % rojo de énfasis (warnbox)
\definecolor{figamber}{HTML}{B8860B}  % ámbar (notebox)
\definecolor{figgray}{HTML}{555555}   % auxiliares, ejes, guías

% ---- CircuiTikz: escala y grosor de los componentes ----
% Guarda \ifdefined: la electrostática (Clases 1-14) no carga `circuitikz`, y
% sin la guarda el \input revienta con `Undefined control sequence \ctikzset`.
% Cuando el curso llegue a circuitos, el bloque se activa solo.
\ifdefined\ctikzset
  \ctikzset{
    bipoles/thickness=1,
    resistors/scale=0.9,
    capacitors/scale=0.9,
  }
\fi

% ---- TikZ: estilos reutilizables ----
% Nota: las puntas se declaran como `-{Latex[...]}` y no como `->`. La forma
% con llaves NO contiene un `>` literal, así que es inmune al `>` activo que
% introduce babel español (ver CLAUDE.md §7.3.3).
\usetikzlibrary{arrows.meta,calc,angles,quotes}

\tikzset{
  figline/.style={figblue, line width=0.9pt},
  figaux/.style={figgray, dashed, line width=0.6pt},
  % vectores
  figvec/.style={figblue, line width=0.9pt, -{Latex[length=2.2mm,width=1.6mm]}},
  figvecres/.style={figred, line width=1.3pt, -{Latex[length=2.6mm,width=1.9mm]}},
  figvecaux/.style={figgray, line width=0.7pt, -{Latex[length=1.8mm,width=1.3mm]}},
  figaxis/.style={figgray, line width=0.6pt, -{Latex[length=2mm,width=1.4mm]}},
  % cargas puntuales
  figcarga/.style={circle, inner sep=0pt, minimum size=5.4pt, draw=none},
  figpos/.style={figcarga, fill=figred},
  figneg/.style={figcarga, fill=figblue},
  figneutra/.style={figcarga, fill=figgray},
  % etiquetas
  figlbl/.style={font=\small},
  figlblaux/.style={font=\footnotesize, figgray},
}

% ---- pgfplots: estilo base de las gráficas ----
\pgfplotsset{
  emfig/.style={
    width=11cm, height=6.5cm,
    axis lines=left,
    xlabel style={font=\small},
    ylabel style={font=\small},
    tick label style={font=\footnotesize},
    every axis plot/.append style={line width=1pt},
    grid=major,
    grid style={figgray!20},
    legend cell align=left,
    legend style={font=\footnotesize, draw=figgray!40},
  },
}


% ---- Title ----
\title{\textbf{Clase 11: Esfera dieléctrica en campo uniforme} \\[0.3em]
        \large{Poisson y Laplace en un dieléctrico, y primeros pasos en energía electrostática}}
\author{Transcripto y expandido --- Curso Electromagnetismo (OpenFING)}
\date{}

\begin{document}

\maketitle
\thispagestyle{empty}
\newpage

\tableofcontents
\newpage

\section{Repaso: medios dieléctricos}

La clase abre con un repaso breve de lo que, según dice el docente, \textbf{les
presentó Julia la semana pasada}: qué cambia en la electrostática cuando hay un
material dieléctrico.

\begin{notebox}[Dieléctrico = aislante]
Las cargas que uno pone en un dieléctrico quedan fijas donde se las puso; no se
desplazan como en un conductor.
\end{notebox}

La dificultad de fondo es determinar \textbf{cómo reacciona el material} a un
campo aplicado: la relación entre la polarización y el campo. Eso es en general
difícil, y lo aborda la \textbf{mecánica estadística} ---la descripción de
sistemas con un número enorme de partículas, en este caso las moléculas del
material---, que no es parte del curso. Aun así hay cosas generales que se pueden
escribir.

\subsection{Lo que sigue valiendo}

Con o sin dieléctrico, el campo sigue derivando de un potencial,
$\vec E = -\grad\phi$. Además se introduce el \textbf{desplazamiento eléctrico}
\[
\vec D = \varepsilon_0 \vec E + \vec P,
\]
que obedece una ley de Gauss con la forma de siempre, pero donde la fuente no es
la carga total:

\begin{keybox}[Ley de Gauss en presencia de dieléctricos]
\[
\boxed{\ \div\vec D = \rho_{\text{libre}}\ },
\qquad
\rho_{\text{total}} = \rho_{\text{libre}} + \rho_{\text{pol}} .
\]
Lo mismo vale para las densidades superficiales.
\end{keybox}

\begin{warnbox}[«Libre» hay que manejarlo con cuidado]
No quiere decir libre de moverse, como en un conductor. Quiere decir
\textbf{externa}: carga a la que un agente externo puede acceder, que puede estar
incrustada y fija en el dieléctrico. Lo que la define es que \textbf{no es carga
de polarización}. En lo que sigue, cuando no se aclara nada, $\rho$ y $\sigma$
son las libres.
\end{warnbox}

Dicho así, esto tiene poca información: la magnitud observable es $\vec E$, y
mientras no se diga cómo se relacionan $\vec E$ y $\vec D$ (o $\vec E$ y
$\vec P$) la ecuación no alcanza para resolver nada.

\subsection{Materiales isótropos, homogéneos y lineales}

La relación puede ser no trivial, pero en casos comunes se simplifica mucho.

\textbf{Isótropo y homogéneo.} Isótropo quiere decir sin direcciones
privilegiadas. Entonces $\vec D$ sólo puede ser paralelo a $\vec E$, con un factor
que depende a lo sumo del módulo:
\[
\vec D = \varepsilon\!\left(|\vec E|\right) \vec E .
\]
Esto supone además un material \textbf{sin polarización espontánea}: sin campo
aplicado no hay polarización.

\textbf{Lineal.} La mayor parte de los dieléctricos cumplen además que ese factor
no depende del campo:

\begin{keybox}[Medio lineal]
\[
\boxed{\ \vec D = \varepsilon \vec E\ },
\qquad \varepsilon = \text{cte},
\qquad
\boxed{\ K = \frac{\varepsilon}{\varepsilon_0}\ }.
\]
$\varepsilon$ es la \textbf{permitividad eléctrica} del material y $K$ la
\textbf{constante dieléctrica}.
\end{keybox}

Sin material no hay polarización y $\vec D = \varepsilon_0\vec E$: el caso lineal
es la misma relación con otra constante, propia de cada material.

\begin{notebox}[Por qué la relación suele ser lineal]
Los campos eléctricos \emph{dentro} de los materiales ---los que ligan electrones
y núcleos--- son enormes, mucho más grandes que cualquier campo que se fabrique en
un laboratorio. A los efectos del material, el campo aplicado es casi cero.
Entonces $\varepsilon(|\vec E|)$ puede verse como el comienzo de un desarrollo de
Taylor alrededor de $E=0$, y reemplazarlo por $\varepsilon(0)$ es una muy buena
aproximación. Para ver la dependencia habría que ir a campos que en la práctica
no se alcanzan.
\end{notebox}

\subsection{Qué pasa si se levanta alguna hipótesis}

\begin{table}[H]
  \centering
  \begin{tabularx}{\textwidth}{lX}
    \toprule
    Hipótesis que falla & Consecuencia \\
    \midrule
    Homogéneo & $\varepsilon$ depende del punto, $\varepsilon(\vec r)$ (p.\,ej.\ un material de densidad variable) \\
    Isótropo (pero lineal y homogéneo) & $\varepsilon$ pasa a ser una \textbf{matriz}, y $\vec D$ ya no es colineal con $\vec E$ \\
    Sin polarización espontánea & Materiales \textbf{ferroeléctricos}: quedan polarizados al apagar el campo \\
    \bottomrule
  \end{tabularx}
\end{table}

\begin{notebox}[Qué materiales cumplen todo]
Los fluidos (gases y líquidos), los sólidos amorfos y algunos cristales. Muchos
sólidos cristalinos no: la estructura del cristal impone direcciones
privilegiadas y los vuelve anisótropos.

Los ferroeléctricos existen ---el docente comenta que en el piso de arriba hay
gente que los estudia--- pero no son tan comunes, y se ignoran en el curso. El
fenómeno análogo en magnetismo, la magnetización espontánea, sí es muy común: son
los imanes.
\end{notebox}

\section{Poisson y Laplace en un dieléctrico}

Con las condiciones de borde que se habían visto la clase anterior, se entra en
problemas electrostáticos con dieléctricos (punto 3.8 de las notas del docente),
siempre con dieléctricos \textbf{homogéneos, isótropos y lineales}. Puede haber
varios, y en cada uno $\vec D = \varepsilon\vec E$ con su propia $\varepsilon$.

Como $\varepsilon$ es constante dentro de cada material, sale de la divergencia:
\[
\div\vec D = \varepsilon\, \div\vec E = \rho .
\]
Es la ley de Gauss del vacío con $\varepsilon$ en lugar de $\varepsilon_0$.
Sustituyendo $\vec E = -\grad\phi$:

\begin{keybox}[Ecuación de Poisson en un dieléctrico]
\[
\boxed{\ \laplacian\phi = -\frac{\rho}{\varepsilon}\ }
\]
válida en cada dieléctrico, cada uno con su constante. Si además no hay carga
libre, $\rho = 0$, vale la \textbf{ecuación de Laplace}, $\laplacian\phi = 0$.
\end{keybox}

\begin{notebox}[Lo interesante]
Laplace vale aunque el medio esté polarizado y haya densidad de carga de
polarización. Lo único que se pide es que la carga \textbf{externa} sea nula y el
material homogéneo, isótropo y lineal.
\end{notebox}

\section{Esfera dieléctrica en un campo uniforme}

Es el ejemplo al que se dedica el grueso de la clase. En una región con un campo
externo \textbf{uniforme} $\vec E_0$ se coloca una esfera de radio $a$, pero no
conductora como la de clases atrás, sino de \textbf{material dieléctrico}
homogéneo, lineal, isótropo, sin polarización espontánea y \textbf{descargado}.

\begin{notebox}[«Descargada» quiere decir sin carga externa]
El material igual se va a polarizar ---sus moléculas responden al campo
aplicado--- y va a aparecer carga de polarización, superficial, volumétrica o
ambas. Cuál de las dos, se ve problema por problema.
\end{notebox}

\subsection{Planteo y las dos zonas}

¿Dónde vale Laplace? \textbf{Afuera} sí: es vacío, sin carga. \textbf{Adentro}
también: no hay carga libre. \textbf{En el borde no}: ahí $\varepsilon$ tiene un
escalón. Si se lo pensara como un único $\varepsilon(\vec r)$, sería homogéneo
afuera y homogéneo adentro, pero no en $r = a$.

\begin{figure}[H]
  \centering
  % Clase 11 — planteo de la esfera dielectrica en campo uniforme.
% Lo que la figura tiene que mostrar: las DOS zonas (1 afuera, 2 adentro), el
% escalon de permitividad en r = a, y las coordenadas (r, theta) medidas desde
% el eje z, que es la direccion de E0. Las condiciones de borde van en la tabla
% del texto, no aca.
\begin{tikzpicture}[scale=1]

  % campo aplicado (lejos, a la izquierda)
  \foreach \yy in {-1.6,-0.8,0,0.8,1.6}{
    \draw[figvec] (-4.6,\yy) -- (-3.5,\yy);}
  \node[figlbl, figblue] at (-4.05,2.15) {$\vec E_0$};

  % esfera
  \fill[figamber!12] (0,0) circle (1.6);
  \draw[figamber, line width=0.9pt] (0,0) circle (1.6);

  % eje z
  \draw[figaxis] (-2.6,0) -- (3.4,0) node[right, figlbl, figgray] {$z$};

  % radio a
  \draw[figvecaux] (0,0) -- (250:1.6);
  \node[figlblaux, fill=figamber!12, inner sep=1pt] at (236:0.95) {$a$};

  % punto de campo y coordenadas
  \coordinate (O) at (0,0);
  \coordinate (P) at (40:2.8);
  \draw[figvec] (O) -- (P);
  \fill[figblue] (P) circle (1.6pt);
  \node[figlbl, figblue, fill=white, inner sep=1pt] at ($(O)!0.62!(P)+(0.05,0.3)$) {$r$};
  \draw[figgray, line width=0.6pt] (0.75,0) arc (0:40:0.75);
  \node[figlblaux] at (20:1.05) {$\theta$};

  % zonas
  \node[figlbl, figamber!80!black] at (-0.55,0.75) {zona 2};
  \node[figlblaux, figamber!80!black] at (-0.55,0.35) {$\varepsilon$};
  \node[figlbl] at (2.55,-1.55) {zona 1};
  \node[figlblaux] at (2.55,-1.95) {$\varepsilon_0$};

\end{tikzpicture}

  \caption{El planteo: una esfera de radio $a$ y permitividad $\varepsilon$ en un
  campo uniforme $\vec E_0$ según $z$. Laplace vale en cada zona por separado,
  pero no en $r=a$, donde la permitividad salta de $\varepsilon$ a
  $\varepsilon_0$. $\theta$ se mide desde la dirección del campo aplicado.}
\end{figure}

Además el problema tiene \textbf{simetría azimutal} en coordenadas esféricas (con
el eje $z$ en la dirección de $\vec E_0$), así que en cada zona se conoce la forma
general de la solución ---la base hallada en clases anteriores---. Se llama
\textbf{zona 1} a la exterior y \textbf{zona 2} a la interior:
\begin{align*}
\phi_1(r,\theta) &= \sum_{n=0}^{\infty}
\left(a_n^{(1)} r^n + \frac{b_n^{(1)}}{r^{n+1}}\right) P_n(\cos\theta),\\
\phi_2(r,\theta) &= \sum_{n=0}^{\infty}
\left(a_n^{(2)} r^n + \frac{b_n^{(2)}}{r^{n+1}}\right) P_n(\cos\theta).
\end{align*}

\begin{notebox}
Los coeficientes de una zona \textbf{no} son los de la otra: para que
coincidieran, las zonas tendrían que estar unidas continuamente por una región
donde valga Laplace. Van a quedar vinculados por las condiciones de borde, pero no
son iguales.
\end{notebox}

\subsection{Las cinco condiciones}

\begin{table}[H]
  \centering
  \begin{tabularx}{\textwidth}{c l c X}
    \toprule
    & Condición & Dónde & De dónde sale \\
    \midrule
    \textbf{A} & $\phi_1 \sim -E_0\, r\cos\theta$ & $r\to\infty$ & Lejos, el campo tiende al aplicado \\
    \textbf{B} & $b_0^{(1)} = 0$ & $r\to\infty$ & La carga neta de polarización es cero \\
    \textbf{C} & $b_n^{(2)} = 0\ \ \forall n$ & $r = 0$ & No hay carga concentrada en el origen \\
    \textbf{D} & $\phi_1(a,\theta) = \phi_2(a,\theta)$ & $r = a$ & El potencial es continuo \\
    \textbf{E} & $D_{r}^{(1)}(a,\theta) = D_{r}^{(2)}(a,\theta)$ & $r = a$ & No hay carga libre superficial \\
    \bottomrule
  \end{tabularx}
\end{table}

\textbf{A.} A gran distancia $\vec E_1 \to \vec E_0 = E_0\hat z$, cuyo potencial
es $-E_0 z = -E_0\, r\cos\theta$. Igual que con la esfera conductora.

\textbf{B.} El efecto de la polarización se puede reemplazar por el de sus
densidades de carga de polarización, que están \textbf{localizadas} en la esfera.
Lejos vale entonces el desarrollo multipolar. El término monopolar es el de la
carga neta, y la carga neta de polarización es \textbf{cero}: el desarrollo
arranca en el dipolo. El término que se comporta como una carga puntual es
$b_0^{(1)}/r$ (recordar que $P_0 = 1$), así que $b_0^{(1)}=0$.

\begin{notebox}[Nada que venga de la esfera puede crecer]
Por el mismo argumento, \textbf{ningún} término que provenga de la esfera puede
crecer con $r$: todo el desarrollo multipolar va en potencias negativas, y el
efecto de una distribución localizada decrece lejos. De ahí que $a_n^{(1)} = 0$
para $n \ge 2$. El único que crece es $a_1^{(1)}$, y ese \textbf{no es de la
esfera}: es del campo aplicado. (Crece el potencial; el campo, no.)
\end{notebox}

\textbf{C.} Un potencial singular en el origen exigiría una carga infinitamente
concentrada ahí. Hay densidad de polarización, pero ninguna carga concentrada.

\textbf{D.} Las discontinuidades de $\vec E$ son \textbf{finitas}: la componente
normal salta cuando hay densidad superficial, pero salta un valor finito. Una
función cuya derivada sólo tiene discontinuidades finitas es continua ---no
derivable, pero continua---. Eso dejaría de valer con cargas concentradas en
líneas o puntos, que acá no hay.

\textbf{E.} La diferencia de las componentes normales de $\vec D$ es proporcional
a la densidad superficial libre, y en $r=a$ no hay: la esfera no tiene carga
externa. Como la superficie es esférica, la componente normal es la radial.

\begin{notebox}[Por qué dos condiciones en el borde]
La ecuación es de \textbf{segundo orden}, y en cada zona hacen falta típicamente
dos.
\end{notebox}

\subsection{Lo que matan A, B y C}

«Hay que armarse de valor»: se resuelve por \textbf{fuerza bruta}, con todos los
términos, antes de ver la versión astuta (§\ref{sec:astuta}).

Con \textbf{A} y \textbf{B}, en la zona exterior sobrevive el término lineal del
campo aplicado, que fija $a_1^{(1)} = -E_0$, y las potencias negativas desde
$n = 1$. Queda además una constante $a_0^{(1)}$ que nada prohíbe; como el
potencial está definido a menos de una constante, \textbf{se elige}
$a_0^{(1)} = 0$. Es una elección de nivel, no física.
\[
\phi_1(r,\theta) = -E_0\, r\cos\theta
+ \sum_{n=1}^{\infty} \frac{b_n^{(1)}}{r^{n+1}} P_n(\cos\theta).
\]
Con \textbf{C}, en la zona interior sólo quedan las potencias positivas:
\[
\phi_2(r,\theta) = \sum_{n=0}^{\infty} a_n^{(2)} r^{n} P_n(\cos\theta).
\]

\begin{warnbox}[Cuidado con tirar también $a_0^{(2)}$]
La tentación es decir «es una constante, la pongo en cero». Pero la constante
aditiva se puede elegir \textbf{una sola vez}: al fijar $a_0^{(1)}=0$, la de
adentro queda determinada por la continuidad. En este problema también va a dar
cero ---«spoiler»---, pero es un accidente del problema, no una regla.
\end{warnbox}

\subsection{Continuidad del potencial}

Se evalúan ambas expresiones en $r=a$ y se igualan para todo $\theta$:
\[
-E_0\, a\cos\theta + \sum_{n=1}^{\infty} \frac{b_n^{(1)}}{a^{n+1}} P_n(\cos\theta)
= \sum_{n=0}^{\infty} a_n^{(2)} a^{n} P_n(\cos\theta).
\]
Los polinomios de Legendre son un conjunto infinito de polinomios, cada uno de
grado distinto: \textbf{linealmente independientes}, una base. Se puede
identificar \textbf{término a término} (como ya se hizo con la esfera conductora):
\begin{itemize}
  \item \textbf{$n=0$.} A la izquierda no hay término de orden cero; a la derecha
    está $a_0^{(2)}$. Entonces $a_0^{(2)} = 0$.
  \item \textbf{$n=1$.} Hay dos contribuciones a la izquierda, porque
    $P_1(\cos\theta) = \cos\theta$ es también el del campo aplicado:
    \[
    -E_0\, a + \frac{b_1^{(1)}}{a^2} = a_1^{(2)}\, a .
    \]
  \item \textbf{$n\ge 2$.} Todas tienen la misma forma:
    \[
    \frac{b_n^{(1)}}{a^{n+1}} = a_n^{(2)} a^n
    \quad\Longrightarrow\quad
    b_n^{(1)} = a_n^{(2)}\, a^{2n+1}.
    \]
\end{itemize}
Los casos especiales son $n=0$ y $n=1$; a partir de $n=2$ son todos iguales.

\begin{notebox}[Todavía no terminó]
Es un sistema lineal \textbf{infinito}, que determina los coeficientes de una
zona en términos de los de la otra. Falta la condición E.
\end{notebox}

\subsection{Continuidad de la componente normal de D}

En cada región $\vec D = \varepsilon\vec E = -\varepsilon\grad\phi$, y la
componente radial del gradiente en esféricas es la derivada respecto de $r$:
\[
D_r^{(1)} = -\varepsilon_0 \pdv{\phi_1}{r},
\qquad
D_r^{(2)} = -\varepsilon\, \pdv{\phi_2}{r}.
\]
Usando $\dv{r}\, r^{-(n+1)} = -(n+1)\, r^{-(n+2)}$ y evaluando en $r=a$, la
condición E se escribe, para todo $\theta$,
\[
\varepsilon_0 E_0\cos\theta
+ \varepsilon_0 \sum_{n=1}^{\infty} \frac{(n+1)\, b_n^{(1)}}{a^{n+2}} P_n(\cos\theta)
= -\varepsilon \sum_{n=1}^{\infty} n\, a_n^{(2)} a^{n-1} P_n(\cos\theta).
\]
(La suma de la derecha arranca en $n=1$ porque ya se sabe que $a_0^{(2)}=0$, y
además su derivada es nula.) Identificando otra vez término a término: no hay
componente $n=0$, y
\begin{align*}
n=1:&\quad \varepsilon_0 E_0 + \frac{2\varepsilon_0\, b_1^{(1)}}{a^3}
= -\varepsilon\, a_1^{(2)},\\
n\ge 2:&\quad \varepsilon_0\,\frac{(n+1)\, b_n^{(1)}}{a^{n+2}}
= -\varepsilon\, n\, a_n^{(2)} a^{n-1}.
\end{align*}

\subsection{Todos los modos con \texorpdfstring{$n \ge 2$}{n ≥ 2} son nulos}

«Muchas ecuaciones, pero casi todas son triviales.» Sustituyendo
$b_n^{(1)} = a_n^{(2)} a^{2n+1}$ de la continuidad en la ecuación de $n\ge2$ y
simplificando el factor $a^{n-1}$:
\[
a_n^{(2)} \left[\varepsilon_0 (n+1) + \varepsilon\, n\right] = 0 .
\]
El corchete es estrictamente positivo, así que
\[
\boxed{\ a_n^{(2)} = b_n^{(1)} = 0 \qquad \forall\, n \ge 2\ }.
\]
Hay «demasiadas ecuaciones» para esas incógnitas: dos condiciones homogéneas por
cada $n$ sólo admiten la solución nula.

\begin{notebox}[Lectura física]
El sistema tiene tanta simetría que la polarización genera \textbf{sólo un
dipolo}. Todos los demás momentos multipolares son cero.
\end{notebox}

\subsection{El sistema para \texorpdfstring{$n = 1$}{n = 1} y la solución}

Quedan dos ecuaciones con dos incógnitas ---ahora \textbf{inhomogéneas}, por los
términos en $E_0$---:
\[
-E_0 + \frac{b_1^{(1)}}{a^3} = a_1^{(2)},
\qquad
\varepsilon_0 E_0 + \frac{2\varepsilon_0\, b_1^{(1)}}{a^3} = -\varepsilon\, a_1^{(2)}.
\]
Reemplazando la primera en la segunda se despeja $b_1^{(1)}$, y con él
$a_1^{(2)}$. (El docente da directamente el resultado final para no equivocarse
en el pizarrón.)
\[
b_1^{(1)} = a^3 E_0\, \frac{\varepsilon-\varepsilon_0}{\varepsilon+2\varepsilon_0},
\qquad
a_1^{(2)} = -\frac{3\varepsilon_0 E_0}{\varepsilon+2\varepsilon_0}.
\]

\begin{keybox}[Esfera dieléctrica en campo uniforme]
En términos de $K = \varepsilon/\varepsilon_0$:
\[
\boxed{\ \phi_1(r,\theta) = -E_0\, r\cos\theta
+ a^3 E_0\, \frac{K-1}{K+2}\, \frac{\cos\theta}{r^2}\ }
\]
\[
\boxed{\ \phi_2(r,\theta) = -\frac{3E_0}{K+2}\, r\cos\theta\ }
\]
\end{keybox}

\begin{notebox}[Qué dice cada uno]
Afuera: el potencial del campo aplicado más el de un \textbf{dipolo} centrado en
la esfera. Comparando con $p\cos\theta/(4\pi\varepsilon_0 r^2)$, el momento es
$p = 4\pi\varepsilon_0\, a^3 E_0\,(K-1)/(K+2)$, en la dirección de $\vec E_0$.
Adentro: un término proporcional a $r\cos\theta = z$.
\end{notebox}

\section{Discusión de la solución}

\subsection{Campo, desplazamiento y polarización adentro}

Como $\phi_2$ es una constante por $z$, conviene pensarlo en cartesianas: el campo
interior es la derivada respecto de $z$ cambiada de signo,
\[
\boxed{\ \vec E_2 = \frac{3}{K+2}\, E_0\, \hat z\ }
\]
\textbf{uniforme}, colineal con el aplicado, pero con otro módulo. De ahí
$\vec D_2 = \varepsilon\vec E_2$ y, usando $\vec D = \varepsilon_0\vec E + \vec P$,
\[
\vec P = (\varepsilon-\varepsilon_0)\, \vec E_2
= 3\varepsilon_0\,\frac{K-1}{K+2}\, E_0\, \hat z .
\]

\begin{notebox}[Chequeo de sentido común]
Si $\varepsilon = \varepsilon_0$ ($K=1$), el problema es uno en que no hay nada:
$\vec E_2 = \vec E_0$, $\vec P = 0$ y el término dipolar de afuera desaparece. Una
esfera con las propiedades dieléctricas del vacío es «transparente» para la
electrostática.
\end{notebox}

\begin{notebox}[Por qué adentro el campo es menor]
Con un campo aplicado, la parte negativa de cada molécula tiende a acercarse a
las cargas positivas que generan el campo y la positiva a las negativas. Las
moléculas vibran, pero en promedio se alinean, y el campo que producen
\textbf{compensa en parte, no del todo}, el original. Como
$\varepsilon \ge \varepsilon_0$ siempre, el efecto es siempre ese: $K \ge 1$ y
$|\vec E_2| \le E_0$.
\end{notebox}

\subsection{Cómo quedan las líneas de campo}

Afuera, lejos, el campo es uniforme y vale $E_0$. Cerca de la esfera es la suma de
un campo constante y un \textbf{campo dipolar}: el mismo tipo de términos que con
la esfera conductora, pero con otros coeficientes. Adentro las líneas son rectas,
paralelas y horizontales, con otra separación porque el módulo es otro.

\begin{figure}[H]
  \centering
  % Clase 11 — lineas de E para la esfera dielectrica en campo uniforme.
% GENERADA: las coordenadas de afuera salen de integrar (RK4) el campo exacto
% E = E0 zhat + beta (3 (zhat.rhat) rhat - zhat)/r^3, beta = (K-1)/(K+2), con
% K = 4, a = 1, E0 = 1. Semillas en z = -3,3 a alturas +-0,2 ... +-2,2 (paso
% 0,4); se integra hasta z = 0 o hasta tocar r = 1 y se espeja z -> -z (el
% campo es simetrico). No retocar los puntos a mano: regenerar.
% Adentro E2 = 3E0/(K+2) = E0/2, asi que las lineas interiores van con el
% DOBLE de separacion que las de afuera lejos (0,8 contra 0,4). El docente las
% dibujo "mas apretadas"; eso vale para D, no para E (ver notes.tex §4.2).
% Las lineas exteriores que tocan la esfera terminan en la carga de
% polarizacion (sigma_p < 0 del lado izquierdo) y renacen del lado derecho.
\usetikzlibrary{decorations.markings}
\begin{tikzpicture}[scale=1.25,
  decoration={markings, mark=at position 0.12 with {\arrow{Latex[length=2mm,width=1.5mm]}}}]

  % esfera
  \fill[figamber!12] (0,0) circle (1);
  \draw[figamber, line width=0.9pt] (0,0) circle (1);

  % lineas exteriores (calculadas)
  \draw[figline, line width=0.7pt, postaction={decorate}] (-3.300,-2.200) -- (-3.240,-2.199) -- (-3.180,-2.199) -- (-3.120,-2.198) -- (-3.060,-2.197) -- (-3.000,-2.196) -- (-2.940,-2.195) -- (-2.880,-2.195) -- (-2.820,-2.194) -- (-2.760,-2.193) -- (-2.700,-2.192) -- (-2.640,-2.191) -- (-2.580,-2.189) -- (-2.520,-2.188) -- (-2.460,-2.187) -- (-2.400,-2.186) -- (-2.340,-2.185) -- (-2.280,-2.183) -- (-2.220,-2.182) -- (-2.160,-2.180) -- (-2.100,-2.179) -- (-2.040,-2.177) -- (-1.980,-2.175) -- (-1.920,-2.173) -- (-1.860,-2.172) -- (-1.800,-2.170) -- (-1.740,-2.168) -- (-1.680,-2.166) -- (-1.620,-2.163) -- (-1.560,-2.161) -- (-1.501,-2.159) -- (-1.441,-2.157) -- (-1.381,-2.154) -- (-1.321,-2.152) -- (-1.261,-2.149) -- (-1.201,-2.147) -- (-1.141,-2.144) -- (-1.081,-2.141) -- (-1.021,-2.139) -- (-0.961,-2.136) -- (-0.901,-2.133) -- (-0.841,-2.131) -- (-0.781,-2.128) -- (-0.721,-2.126) -- (-0.661,-2.123) -- (-0.601,-2.121) -- (-0.541,-2.119) -- (-0.481,-2.116) -- (-0.421,-2.115) -- (-0.361,-2.113) -- (-0.301,-2.111) -- (-0.242,-2.110) -- (-0.182,-2.109) -- (-0.122,-2.108) -- (-0.062,-2.108) -- (-0.002,-2.108) -- (0.008,-2.108) -- (0.002,-2.108) -- (0.062,-2.108) -- (0.122,-2.108) -- (0.182,-2.109) -- (0.242,-2.110) -- (0.301,-2.111) -- (0.361,-2.113) -- (0.421,-2.115) -- (0.481,-2.116) -- (0.541,-2.119) -- (0.601,-2.121) -- (0.661,-2.123) -- (0.721,-2.126) -- (0.781,-2.128) -- (0.841,-2.131) -- (0.901,-2.133) -- (0.961,-2.136) -- (1.021,-2.139) -- (1.081,-2.141) -- (1.141,-2.144) -- (1.201,-2.147) -- (1.261,-2.149) -- (1.321,-2.152) -- (1.381,-2.154) -- (1.441,-2.157) -- (1.501,-2.159) -- (1.560,-2.161) -- (1.620,-2.163) -- (1.680,-2.166) -- (1.740,-2.168) -- (1.800,-2.170) -- (1.860,-2.172) -- (1.920,-2.173) -- (1.980,-2.175) -- (2.040,-2.177) -- (2.100,-2.179) -- (2.160,-2.180) -- (2.220,-2.182) -- (2.280,-2.183) -- (2.340,-2.185) -- (2.400,-2.186) -- (2.460,-2.187) -- (2.520,-2.188) -- (2.580,-2.189) -- (2.640,-2.191) -- (2.700,-2.192) -- (2.760,-2.193) -- (2.820,-2.194) -- (2.880,-2.195) -- (2.940,-2.195) -- (3.000,-2.196) -- (3.060,-2.197) -- (3.120,-2.198) -- (3.180,-2.199) -- (3.240,-2.199) -- (3.300,-2.200);
  \draw[figline, line width=0.7pt, postaction={decorate}] (-3.300,-1.800) -- (-3.240,-1.799) -- (-3.180,-1.799) -- (-3.120,-1.798) -- (-3.060,-1.797) -- (-3.000,-1.796) -- (-2.940,-1.795) -- (-2.880,-1.794) -- (-2.820,-1.793) -- (-2.760,-1.792) -- (-2.700,-1.791) -- (-2.640,-1.789) -- (-2.580,-1.788) -- (-2.520,-1.787) -- (-2.460,-1.785) -- (-2.400,-1.784) -- (-2.340,-1.782) -- (-2.280,-1.780) -- (-2.220,-1.779) -- (-2.160,-1.777) -- (-2.100,-1.775) -- (-2.040,-1.773) -- (-1.980,-1.770) -- (-1.920,-1.768) -- (-1.860,-1.765) -- (-1.801,-1.763) -- (-1.741,-1.760) -- (-1.681,-1.757) -- (-1.621,-1.754) -- (-1.561,-1.750) -- (-1.501,-1.747) -- (-1.441,-1.743) -- (-1.381,-1.739) -- (-1.321,-1.735) -- (-1.261,-1.731) -- (-1.202,-1.726) -- (-1.142,-1.722) -- (-1.082,-1.717) -- (-1.022,-1.712) -- (-0.962,-1.707) -- (-0.903,-1.702) -- (-0.843,-1.696) -- (-0.783,-1.691) -- (-0.723,-1.685) -- (-0.664,-1.680) -- (-0.604,-1.674) -- (-0.544,-1.669) -- (-0.484,-1.664) -- (-0.425,-1.659) -- (-0.365,-1.655) -- (-0.305,-1.651) -- (-0.245,-1.647) -- (-0.185,-1.644) -- (-0.125,-1.642) -- (-0.065,-1.641) -- (-0.005,-1.641) -- (0.005,-1.641) -- (0.005,-1.641) -- (0.065,-1.641) -- (0.125,-1.642) -- (0.185,-1.644) -- (0.245,-1.647) -- (0.305,-1.651) -- (0.365,-1.655) -- (0.425,-1.659) -- (0.484,-1.664) -- (0.544,-1.669) -- (0.604,-1.674) -- (0.664,-1.680) -- (0.723,-1.685) -- (0.783,-1.691) -- (0.843,-1.696) -- (0.903,-1.702) -- (0.962,-1.707) -- (1.022,-1.712) -- (1.082,-1.717) -- (1.142,-1.722) -- (1.202,-1.726) -- (1.261,-1.731) -- (1.321,-1.735) -- (1.381,-1.739) -- (1.441,-1.743) -- (1.501,-1.747) -- (1.561,-1.750) -- (1.621,-1.754) -- (1.681,-1.757) -- (1.741,-1.760) -- (1.801,-1.763) -- (1.860,-1.765) -- (1.920,-1.768) -- (1.980,-1.770) -- (2.040,-1.773) -- (2.100,-1.775) -- (2.160,-1.777) -- (2.220,-1.779) -- (2.280,-1.780) -- (2.340,-1.782) -- (2.400,-1.784) -- (2.460,-1.785) -- (2.520,-1.787) -- (2.580,-1.788) -- (2.640,-1.789) -- (2.700,-1.791) -- (2.760,-1.792) -- (2.820,-1.793) -- (2.880,-1.794) -- (2.940,-1.795) -- (3.000,-1.796) -- (3.060,-1.797) -- (3.120,-1.798) -- (3.180,-1.799) -- (3.240,-1.799) -- (3.300,-1.800);
  \draw[figline, line width=0.7pt, postaction={decorate}] (-3.300,-1.400) -- (-3.240,-1.399) -- (-3.180,-1.399) -- (-3.120,-1.398) -- (-3.060,-1.397) -- (-3.000,-1.396) -- (-2.940,-1.395) -- (-2.880,-1.394) -- (-2.820,-1.393) -- (-2.760,-1.392) -- (-2.700,-1.391) -- (-2.640,-1.389) -- (-2.580,-1.388) -- (-2.520,-1.386) -- (-2.460,-1.385) -- (-2.400,-1.383) -- (-2.340,-1.381) -- (-2.280,-1.379) -- (-2.220,-1.377) -- (-2.160,-1.375) -- (-2.100,-1.372) -- (-2.040,-1.370) -- (-1.980,-1.367) -- (-1.921,-1.364) -- (-1.861,-1.361) -- (-1.801,-1.357) -- (-1.741,-1.353) -- (-1.681,-1.349) -- (-1.621,-1.345) -- (-1.561,-1.340) -- (-1.502,-1.335) -- (-1.442,-1.330) -- (-1.382,-1.324) -- (-1.322,-1.317) -- (-1.263,-1.311) -- (-1.203,-1.303) -- (-1.144,-1.295) -- (-1.084,-1.287) -- (-1.025,-1.277) -- (-0.966,-1.267) -- (-0.907,-1.256) -- (-0.848,-1.245) -- (-0.789,-1.232) -- (-0.731,-1.219) -- (-0.673,-1.204) -- (-0.615,-1.189) -- (-0.557,-1.172) -- (-0.500,-1.155) -- (-0.443,-1.136) -- (-0.386,-1.116) -- (-0.330,-1.096) -- (-0.273,-1.075) -- (-0.217,-1.054) -- (-0.161,-1.034) -- (-0.103,-1.017) -- (-0.044,-1.005) -- (0.006,-1.003) -- (0.044,-1.005) -- (0.103,-1.017) -- (0.161,-1.034) -- (0.217,-1.054) -- (0.273,-1.075) -- (0.330,-1.096) -- (0.386,-1.116) -- (0.443,-1.136) -- (0.500,-1.155) -- (0.557,-1.172) -- (0.615,-1.189) -- (0.673,-1.204) -- (0.731,-1.219) -- (0.789,-1.232) -- (0.848,-1.245) -- (0.907,-1.256) -- (0.966,-1.267) -- (1.025,-1.277) -- (1.084,-1.287) -- (1.144,-1.295) -- (1.203,-1.303) -- (1.263,-1.311) -- (1.322,-1.317) -- (1.382,-1.324) -- (1.442,-1.330) -- (1.502,-1.335) -- (1.561,-1.340) -- (1.621,-1.345) -- (1.681,-1.349) -- (1.741,-1.353) -- (1.801,-1.357) -- (1.861,-1.361) -- (1.921,-1.364) -- (1.980,-1.367) -- (2.040,-1.370) -- (2.100,-1.372) -- (2.160,-1.375) -- (2.220,-1.377) -- (2.280,-1.379) -- (2.340,-1.381) -- (2.400,-1.383) -- (2.460,-1.385) -- (2.520,-1.386) -- (2.580,-1.388) -- (2.640,-1.389) -- (2.700,-1.391) -- (2.760,-1.392) -- (2.820,-1.393) -- (2.880,-1.394) -- (2.940,-1.395) -- (3.000,-1.396) -- (3.060,-1.397) -- (3.120,-1.398) -- (3.180,-1.399) -- (3.240,-1.399) -- (3.300,-1.400);
  \draw[figline, line width=0.7pt, postaction={decorate}] (-3.300,-1.000) -- (-3.240,-0.999) -- (-3.180,-0.999) -- (-3.120,-0.998) -- (-3.060,-0.997) -- (-3.000,-0.996) -- (-2.940,-0.996) -- (-2.880,-0.995) -- (-2.820,-0.994) -- (-2.760,-0.993) -- (-2.700,-0.992) -- (-2.640,-0.990) -- (-2.580,-0.989) -- (-2.520,-0.988) -- (-2.460,-0.986) -- (-2.400,-0.984) -- (-2.340,-0.983) -- (-2.280,-0.981) -- (-2.220,-0.978) -- (-2.160,-0.976) -- (-2.100,-0.974) -- (-2.040,-0.971) -- (-1.980,-0.968) -- (-1.921,-0.965) -- (-1.861,-0.961) -- (-1.801,-0.957) -- (-1.741,-0.953) -- (-1.681,-0.948) -- (-1.621,-0.943) -- (-1.562,-0.938) -- (-1.502,-0.931) -- (-1.442,-0.924) -- (-1.383,-0.917) -- (-1.323,-0.908) -- (-1.264,-0.899) -- (-1.205,-0.888) -- (-1.146,-0.877) -- (-1.088,-0.864) -- (-1.029,-0.849) -- (-0.972,-0.833) -- (-0.915,-0.814) -- (-0.858,-0.794) -- (-0.803,-0.770) -- (-0.749,-0.744) -- (-0.698,-0.716) -- (-0.698,-0.716);
  \draw[figline, line width=0.7pt] (0.698,-0.716) -- (0.749,-0.744) -- (0.803,-0.770) -- (0.858,-0.794) -- (0.915,-0.814) -- (0.972,-0.833) -- (1.029,-0.849) -- (1.088,-0.864) -- (1.146,-0.877) -- (1.205,-0.888) -- (1.264,-0.899) -- (1.323,-0.908) -- (1.383,-0.917) -- (1.442,-0.924) -- (1.502,-0.931) -- (1.562,-0.938) -- (1.621,-0.943) -- (1.681,-0.948) -- (1.741,-0.953) -- (1.801,-0.957) -- (1.861,-0.961) -- (1.921,-0.965) -- (1.980,-0.968) -- (2.040,-0.971) -- (2.100,-0.974) -- (2.160,-0.976) -- (2.220,-0.978) -- (2.280,-0.981) -- (2.340,-0.983) -- (2.400,-0.984) -- (2.460,-0.986) -- (2.520,-0.988) -- (2.580,-0.989) -- (2.640,-0.990) -- (2.700,-0.992) -- (2.760,-0.993) -- (2.820,-0.994) -- (2.880,-0.995) -- (2.940,-0.996) -- (3.000,-0.996) -- (3.060,-0.997) -- (3.120,-0.998) -- (3.180,-0.999) -- (3.240,-0.999) -- (3.300,-1.000);
  \draw[figline, line width=0.7pt, postaction={decorate}] (-3.300,-0.600) -- (-3.240,-0.600) -- (-3.180,-0.599) -- (-3.120,-0.599) -- (-3.060,-0.598) -- (-3.000,-0.598) -- (-2.940,-0.597) -- (-2.880,-0.596) -- (-2.820,-0.596) -- (-2.760,-0.595) -- (-2.700,-0.594) -- (-2.640,-0.593) -- (-2.580,-0.592) -- (-2.520,-0.591) -- (-2.460,-0.590) -- (-2.400,-0.589) -- (-2.340,-0.587) -- (-2.280,-0.586) -- (-2.220,-0.584) -- (-2.160,-0.582) -- (-2.100,-0.580) -- (-2.040,-0.578) -- (-1.980,-0.576) -- (-1.920,-0.573) -- (-1.860,-0.570) -- (-1.800,-0.567) -- (-1.741,-0.564) -- (-1.681,-0.560) -- (-1.621,-0.555) -- (-1.561,-0.550) -- (-1.501,-0.545) -- (-1.442,-0.539) -- (-1.382,-0.532) -- (-1.323,-0.524) -- (-1.263,-0.515) -- (-1.204,-0.505) -- (-1.145,-0.494) -- (-1.086,-0.481) -- (-1.028,-0.467) -- (-0.970,-0.451) -- (-0.913,-0.433) -- (-0.902,-0.431);
  \draw[figline, line width=0.7pt] (0.902,-0.431) -- (0.913,-0.433) -- (0.970,-0.451) -- (1.028,-0.467) -- (1.086,-0.481) -- (1.145,-0.494) -- (1.204,-0.505) -- (1.263,-0.515) -- (1.323,-0.524) -- (1.382,-0.532) -- (1.442,-0.539) -- (1.501,-0.545) -- (1.561,-0.550) -- (1.621,-0.555) -- (1.681,-0.560) -- (1.741,-0.564) -- (1.800,-0.567) -- (1.860,-0.570) -- (1.920,-0.573) -- (1.980,-0.576) -- (2.040,-0.578) -- (2.100,-0.580) -- (2.160,-0.582) -- (2.220,-0.584) -- (2.280,-0.586) -- (2.340,-0.587) -- (2.400,-0.589) -- (2.460,-0.590) -- (2.520,-0.591) -- (2.580,-0.592) -- (2.640,-0.593) -- (2.700,-0.594) -- (2.760,-0.595) -- (2.820,-0.596) -- (2.880,-0.596) -- (2.940,-0.597) -- (3.000,-0.598) -- (3.060,-0.598) -- (3.120,-0.599) -- (3.180,-0.599) -- (3.240,-0.600) -- (3.300,-0.600);
  \draw[figline, line width=0.7pt, postaction={decorate}] (-3.300,-0.200) -- (-3.240,-0.200) -- (-3.180,-0.200) -- (-3.120,-0.200) -- (-3.060,-0.199) -- (-3.000,-0.199) -- (-2.940,-0.199) -- (-2.880,-0.199) -- (-2.820,-0.198) -- (-2.760,-0.198) -- (-2.700,-0.198) -- (-2.640,-0.197) -- (-2.580,-0.197) -- (-2.520,-0.197) -- (-2.460,-0.196) -- (-2.400,-0.196) -- (-2.340,-0.195) -- (-2.280,-0.195) -- (-2.220,-0.194) -- (-2.160,-0.193) -- (-2.100,-0.193) -- (-2.040,-0.192) -- (-1.980,-0.191) -- (-1.920,-0.190) -- (-1.860,-0.189) -- (-1.800,-0.188) -- (-1.740,-0.186) -- (-1.680,-0.185) -- (-1.620,-0.183) -- (-1.560,-0.181) -- (-1.500,-0.179) -- (-1.440,-0.176) -- (-1.380,-0.173) -- (-1.320,-0.170) -- (-1.261,-0.166) -- (-1.201,-0.162) -- (-1.141,-0.158) -- (-1.081,-0.152) -- (-1.021,-0.147) -- (-0.990,-0.144);
  \draw[figline, line width=0.7pt] (0.990,-0.144) -- (1.021,-0.147) -- (1.081,-0.152) -- (1.141,-0.158) -- (1.201,-0.162) -- (1.261,-0.166) -- (1.320,-0.170) -- (1.380,-0.173) -- (1.440,-0.176) -- (1.500,-0.179) -- (1.560,-0.181) -- (1.620,-0.183) -- (1.680,-0.185) -- (1.740,-0.186) -- (1.800,-0.188) -- (1.860,-0.189) -- (1.920,-0.190) -- (1.980,-0.191) -- (2.040,-0.192) -- (2.100,-0.193) -- (2.160,-0.193) -- (2.220,-0.194) -- (2.280,-0.195) -- (2.340,-0.195) -- (2.400,-0.196) -- (2.460,-0.196) -- (2.520,-0.197) -- (2.580,-0.197) -- (2.640,-0.197) -- (2.700,-0.198) -- (2.760,-0.198) -- (2.820,-0.198) -- (2.880,-0.199) -- (2.940,-0.199) -- (3.000,-0.199) -- (3.060,-0.199) -- (3.120,-0.200) -- (3.180,-0.200) -- (3.240,-0.200) -- (3.300,-0.200);
  \draw[figline, line width=0.7pt, postaction={decorate}] (-3.300,0.200) -- (-3.240,0.200) -- (-3.180,0.200) -- (-3.120,0.200) -- (-3.060,0.199) -- (-3.000,0.199) -- (-2.940,0.199) -- (-2.880,0.199) -- (-2.820,0.198) -- (-2.760,0.198) -- (-2.700,0.198) -- (-2.640,0.197) -- (-2.580,0.197) -- (-2.520,0.197) -- (-2.460,0.196) -- (-2.400,0.196) -- (-2.340,0.195) -- (-2.280,0.195) -- (-2.220,0.194) -- (-2.160,0.193) -- (-2.100,0.193) -- (-2.040,0.192) -- (-1.980,0.191) -- (-1.920,0.190) -- (-1.860,0.189) -- (-1.800,0.188) -- (-1.740,0.186) -- (-1.680,0.185) -- (-1.620,0.183) -- (-1.560,0.181) -- (-1.500,0.179) -- (-1.440,0.176) -- (-1.380,0.173) -- (-1.320,0.170) -- (-1.261,0.166) -- (-1.201,0.162) -- (-1.141,0.158) -- (-1.081,0.152) -- (-1.021,0.147) -- (-0.990,0.144);
  \draw[figline, line width=0.7pt] (0.990,0.144) -- (1.021,0.147) -- (1.081,0.152) -- (1.141,0.158) -- (1.201,0.162) -- (1.261,0.166) -- (1.320,0.170) -- (1.380,0.173) -- (1.440,0.176) -- (1.500,0.179) -- (1.560,0.181) -- (1.620,0.183) -- (1.680,0.185) -- (1.740,0.186) -- (1.800,0.188) -- (1.860,0.189) -- (1.920,0.190) -- (1.980,0.191) -- (2.040,0.192) -- (2.100,0.193) -- (2.160,0.193) -- (2.220,0.194) -- (2.280,0.195) -- (2.340,0.195) -- (2.400,0.196) -- (2.460,0.196) -- (2.520,0.197) -- (2.580,0.197) -- (2.640,0.197) -- (2.700,0.198) -- (2.760,0.198) -- (2.820,0.198) -- (2.880,0.199) -- (2.940,0.199) -- (3.000,0.199) -- (3.060,0.199) -- (3.120,0.200) -- (3.180,0.200) -- (3.240,0.200) -- (3.300,0.200);
  \draw[figline, line width=0.7pt, postaction={decorate}] (-3.300,0.600) -- (-3.240,0.600) -- (-3.180,0.599) -- (-3.120,0.599) -- (-3.060,0.598) -- (-3.000,0.598) -- (-2.940,0.597) -- (-2.880,0.596) -- (-2.820,0.596) -- (-2.760,0.595) -- (-2.700,0.594) -- (-2.640,0.593) -- (-2.580,0.592) -- (-2.520,0.591) -- (-2.460,0.590) -- (-2.400,0.589) -- (-2.340,0.587) -- (-2.280,0.586) -- (-2.220,0.584) -- (-2.160,0.582) -- (-2.100,0.580) -- (-2.040,0.578) -- (-1.980,0.576) -- (-1.920,0.573) -- (-1.860,0.570) -- (-1.800,0.567) -- (-1.741,0.564) -- (-1.681,0.560) -- (-1.621,0.555) -- (-1.561,0.550) -- (-1.501,0.545) -- (-1.442,0.539) -- (-1.382,0.532) -- (-1.323,0.524) -- (-1.263,0.515) -- (-1.204,0.505) -- (-1.145,0.494) -- (-1.086,0.481) -- (-1.028,0.467) -- (-0.970,0.451) -- (-0.913,0.433) -- (-0.902,0.431);
  \draw[figline, line width=0.7pt] (0.902,0.431) -- (0.913,0.433) -- (0.970,0.451) -- (1.028,0.467) -- (1.086,0.481) -- (1.145,0.494) -- (1.204,0.505) -- (1.263,0.515) -- (1.323,0.524) -- (1.382,0.532) -- (1.442,0.539) -- (1.501,0.545) -- (1.561,0.550) -- (1.621,0.555) -- (1.681,0.560) -- (1.741,0.564) -- (1.800,0.567) -- (1.860,0.570) -- (1.920,0.573) -- (1.980,0.576) -- (2.040,0.578) -- (2.100,0.580) -- (2.160,0.582) -- (2.220,0.584) -- (2.280,0.586) -- (2.340,0.587) -- (2.400,0.589) -- (2.460,0.590) -- (2.520,0.591) -- (2.580,0.592) -- (2.640,0.593) -- (2.700,0.594) -- (2.760,0.595) -- (2.820,0.596) -- (2.880,0.596) -- (2.940,0.597) -- (3.000,0.598) -- (3.060,0.598) -- (3.120,0.599) -- (3.180,0.599) -- (3.240,0.600) -- (3.300,0.600);
  \draw[figline, line width=0.7pt, postaction={decorate}] (-3.300,1.000) -- (-3.240,0.999) -- (-3.180,0.999) -- (-3.120,0.998) -- (-3.060,0.997) -- (-3.000,0.996) -- (-2.940,0.996) -- (-2.880,0.995) -- (-2.820,0.994) -- (-2.760,0.993) -- (-2.700,0.992) -- (-2.640,0.990) -- (-2.580,0.989) -- (-2.520,0.988) -- (-2.460,0.986) -- (-2.400,0.984) -- (-2.340,0.983) -- (-2.280,0.981) -- (-2.220,0.978) -- (-2.160,0.976) -- (-2.100,0.974) -- (-2.040,0.971) -- (-1.980,0.968) -- (-1.921,0.965) -- (-1.861,0.961) -- (-1.801,0.957) -- (-1.741,0.953) -- (-1.681,0.948) -- (-1.621,0.943) -- (-1.562,0.938) -- (-1.502,0.931) -- (-1.442,0.924) -- (-1.383,0.917) -- (-1.323,0.908) -- (-1.264,0.899) -- (-1.205,0.888) -- (-1.146,0.877) -- (-1.088,0.864) -- (-1.029,0.849) -- (-0.972,0.833) -- (-0.915,0.814) -- (-0.858,0.794) -- (-0.803,0.770) -- (-0.749,0.744) -- (-0.698,0.716) -- (-0.698,0.716);
  \draw[figline, line width=0.7pt] (0.698,0.716) -- (0.749,0.744) -- (0.803,0.770) -- (0.858,0.794) -- (0.915,0.814) -- (0.972,0.833) -- (1.029,0.849) -- (1.088,0.864) -- (1.146,0.877) -- (1.205,0.888) -- (1.264,0.899) -- (1.323,0.908) -- (1.383,0.917) -- (1.442,0.924) -- (1.502,0.931) -- (1.562,0.938) -- (1.621,0.943) -- (1.681,0.948) -- (1.741,0.953) -- (1.801,0.957) -- (1.861,0.961) -- (1.921,0.965) -- (1.980,0.968) -- (2.040,0.971) -- (2.100,0.974) -- (2.160,0.976) -- (2.220,0.978) -- (2.280,0.981) -- (2.340,0.983) -- (2.400,0.984) -- (2.460,0.986) -- (2.520,0.988) -- (2.580,0.989) -- (2.640,0.990) -- (2.700,0.992) -- (2.760,0.993) -- (2.820,0.994) -- (2.880,0.995) -- (2.940,0.996) -- (3.000,0.996) -- (3.060,0.997) -- (3.120,0.998) -- (3.180,0.999) -- (3.240,0.999) -- (3.300,1.000);
  \draw[figline, line width=0.7pt, postaction={decorate}] (-3.300,1.400) -- (-3.240,1.399) -- (-3.180,1.399) -- (-3.120,1.398) -- (-3.060,1.397) -- (-3.000,1.396) -- (-2.940,1.395) -- (-2.880,1.394) -- (-2.820,1.393) -- (-2.760,1.392) -- (-2.700,1.391) -- (-2.640,1.389) -- (-2.580,1.388) -- (-2.520,1.386) -- (-2.460,1.385) -- (-2.400,1.383) -- (-2.340,1.381) -- (-2.280,1.379) -- (-2.220,1.377) -- (-2.160,1.375) -- (-2.100,1.372) -- (-2.040,1.370) -- (-1.980,1.367) -- (-1.921,1.364) -- (-1.861,1.361) -- (-1.801,1.357) -- (-1.741,1.353) -- (-1.681,1.349) -- (-1.621,1.345) -- (-1.561,1.340) -- (-1.502,1.335) -- (-1.442,1.330) -- (-1.382,1.324) -- (-1.322,1.317) -- (-1.263,1.311) -- (-1.203,1.303) -- (-1.144,1.295) -- (-1.084,1.287) -- (-1.025,1.277) -- (-0.966,1.267) -- (-0.907,1.256) -- (-0.848,1.245) -- (-0.789,1.232) -- (-0.731,1.219) -- (-0.673,1.204) -- (-0.615,1.189) -- (-0.557,1.172) -- (-0.500,1.155) -- (-0.443,1.136) -- (-0.386,1.116) -- (-0.330,1.096) -- (-0.273,1.075) -- (-0.217,1.054) -- (-0.161,1.034) -- (-0.103,1.017) -- (-0.044,1.005) -- (0.006,1.003) -- (0.044,1.005) -- (0.103,1.017) -- (0.161,1.034) -- (0.217,1.054) -- (0.273,1.075) -- (0.330,1.096) -- (0.386,1.116) -- (0.443,1.136) -- (0.500,1.155) -- (0.557,1.172) -- (0.615,1.189) -- (0.673,1.204) -- (0.731,1.219) -- (0.789,1.232) -- (0.848,1.245) -- (0.907,1.256) -- (0.966,1.267) -- (1.025,1.277) -- (1.084,1.287) -- (1.144,1.295) -- (1.203,1.303) -- (1.263,1.311) -- (1.322,1.317) -- (1.382,1.324) -- (1.442,1.330) -- (1.502,1.335) -- (1.561,1.340) -- (1.621,1.345) -- (1.681,1.349) -- (1.741,1.353) -- (1.801,1.357) -- (1.861,1.361) -- (1.921,1.364) -- (1.980,1.367) -- (2.040,1.370) -- (2.100,1.372) -- (2.160,1.375) -- (2.220,1.377) -- (2.280,1.379) -- (2.340,1.381) -- (2.400,1.383) -- (2.460,1.385) -- (2.520,1.386) -- (2.580,1.388) -- (2.640,1.389) -- (2.700,1.391) -- (2.760,1.392) -- (2.820,1.393) -- (2.880,1.394) -- (2.940,1.395) -- (3.000,1.396) -- (3.060,1.397) -- (3.120,1.398) -- (3.180,1.399) -- (3.240,1.399) -- (3.300,1.400);
  \draw[figline, line width=0.7pt, postaction={decorate}] (-3.300,1.800) -- (-3.240,1.799) -- (-3.180,1.799) -- (-3.120,1.798) -- (-3.060,1.797) -- (-3.000,1.796) -- (-2.940,1.795) -- (-2.880,1.794) -- (-2.820,1.793) -- (-2.760,1.792) -- (-2.700,1.791) -- (-2.640,1.789) -- (-2.580,1.788) -- (-2.520,1.787) -- (-2.460,1.785) -- (-2.400,1.784) -- (-2.340,1.782) -- (-2.280,1.780) -- (-2.220,1.779) -- (-2.160,1.777) -- (-2.100,1.775) -- (-2.040,1.773) -- (-1.980,1.770) -- (-1.920,1.768) -- (-1.860,1.765) -- (-1.801,1.763) -- (-1.741,1.760) -- (-1.681,1.757) -- (-1.621,1.754) -- (-1.561,1.750) -- (-1.501,1.747) -- (-1.441,1.743) -- (-1.381,1.739) -- (-1.321,1.735) -- (-1.261,1.731) -- (-1.202,1.726) -- (-1.142,1.722) -- (-1.082,1.717) -- (-1.022,1.712) -- (-0.962,1.707) -- (-0.903,1.702) -- (-0.843,1.696) -- (-0.783,1.691) -- (-0.723,1.685) -- (-0.664,1.680) -- (-0.604,1.674) -- (-0.544,1.669) -- (-0.484,1.664) -- (-0.425,1.659) -- (-0.365,1.655) -- (-0.305,1.651) -- (-0.245,1.647) -- (-0.185,1.644) -- (-0.125,1.642) -- (-0.065,1.641) -- (-0.005,1.641) -- (0.005,1.641) -- (0.005,1.641) -- (0.065,1.641) -- (0.125,1.642) -- (0.185,1.644) -- (0.245,1.647) -- (0.305,1.651) -- (0.365,1.655) -- (0.425,1.659) -- (0.484,1.664) -- (0.544,1.669) -- (0.604,1.674) -- (0.664,1.680) -- (0.723,1.685) -- (0.783,1.691) -- (0.843,1.696) -- (0.903,1.702) -- (0.962,1.707) -- (1.022,1.712) -- (1.082,1.717) -- (1.142,1.722) -- (1.202,1.726) -- (1.261,1.731) -- (1.321,1.735) -- (1.381,1.739) -- (1.441,1.743) -- (1.501,1.747) -- (1.561,1.750) -- (1.621,1.754) -- (1.681,1.757) -- (1.741,1.760) -- (1.801,1.763) -- (1.860,1.765) -- (1.920,1.768) -- (1.980,1.770) -- (2.040,1.773) -- (2.100,1.775) -- (2.160,1.777) -- (2.220,1.779) -- (2.280,1.780) -- (2.340,1.782) -- (2.400,1.784) -- (2.460,1.785) -- (2.520,1.787) -- (2.580,1.788) -- (2.640,1.789) -- (2.700,1.791) -- (2.760,1.792) -- (2.820,1.793) -- (2.880,1.794) -- (2.940,1.795) -- (3.000,1.796) -- (3.060,1.797) -- (3.120,1.798) -- (3.180,1.799) -- (3.240,1.799) -- (3.300,1.800);
  \draw[figline, line width=0.7pt, postaction={decorate}] (-3.300,2.200) -- (-3.240,2.199) -- (-3.180,2.199) -- (-3.120,2.198) -- (-3.060,2.197) -- (-3.000,2.196) -- (-2.940,2.195) -- (-2.880,2.195) -- (-2.820,2.194) -- (-2.760,2.193) -- (-2.700,2.192) -- (-2.640,2.191) -- (-2.580,2.189) -- (-2.520,2.188) -- (-2.460,2.187) -- (-2.400,2.186) -- (-2.340,2.185) -- (-2.280,2.183) -- (-2.220,2.182) -- (-2.160,2.180) -- (-2.100,2.179) -- (-2.040,2.177) -- (-1.980,2.175) -- (-1.920,2.173) -- (-1.860,2.172) -- (-1.800,2.170) -- (-1.740,2.168) -- (-1.680,2.166) -- (-1.620,2.163) -- (-1.560,2.161) -- (-1.501,2.159) -- (-1.441,2.157) -- (-1.381,2.154) -- (-1.321,2.152) -- (-1.261,2.149) -- (-1.201,2.147) -- (-1.141,2.144) -- (-1.081,2.141) -- (-1.021,2.139) -- (-0.961,2.136) -- (-0.901,2.133) -- (-0.841,2.131) -- (-0.781,2.128) -- (-0.721,2.126) -- (-0.661,2.123) -- (-0.601,2.121) -- (-0.541,2.119) -- (-0.481,2.116) -- (-0.421,2.115) -- (-0.361,2.113) -- (-0.301,2.111) -- (-0.242,2.110) -- (-0.182,2.109) -- (-0.122,2.108) -- (-0.062,2.108) -- (-0.002,2.108) -- (0.008,2.108) -- (0.002,2.108) -- (0.062,2.108) -- (0.122,2.108) -- (0.182,2.109) -- (0.242,2.110) -- (0.301,2.111) -- (0.361,2.113) -- (0.421,2.115) -- (0.481,2.116) -- (0.541,2.119) -- (0.601,2.121) -- (0.661,2.123) -- (0.721,2.126) -- (0.781,2.128) -- (0.841,2.131) -- (0.901,2.133) -- (0.961,2.136) -- (1.021,2.139) -- (1.081,2.141) -- (1.141,2.144) -- (1.201,2.147) -- (1.261,2.149) -- (1.321,2.152) -- (1.381,2.154) -- (1.441,2.157) -- (1.501,2.159) -- (1.560,2.161) -- (1.620,2.163) -- (1.680,2.166) -- (1.740,2.168) -- (1.800,2.170) -- (1.860,2.172) -- (1.920,2.173) -- (1.980,2.175) -- (2.040,2.177) -- (2.100,2.179) -- (2.160,2.180) -- (2.220,2.182) -- (2.280,2.183) -- (2.340,2.185) -- (2.400,2.186) -- (2.460,2.187) -- (2.520,2.188) -- (2.580,2.189) -- (2.640,2.191) -- (2.700,2.192) -- (2.760,2.193) -- (2.820,2.194) -- (2.880,2.195) -- (2.940,2.195) -- (3.000,2.196) -- (3.060,2.197) -- (3.120,2.198) -- (3.180,2.199) -- (3.240,2.199) -- (3.300,2.200);

  % lineas interiores: uniformes, E2 = E0/2
  \foreach \xx in {-0.8,0,0.8}{
    \pgfmathsetmacro{\zz}{sqrt(1-\xx*\xx)}
    \draw[figline, line width=0.7pt, -{Latex[length=2mm,width=1.5mm]}] (-\zz,\xx) -- (\zz,\xx);}

  % rotulos
  \node[figlbl, figblue, fill=white, inner sep=1.5pt] at (-3.0,2.62) {$\vec E_0$};
  \node[figlbl, figamber!80!black, fill=figamber!12, inner sep=1pt] at (0,0.4) {$K$};
  \node[figlblaux, fill=white, inner sep=1pt] at (2.75,-2.45) {$\varepsilon_0$};

\end{tikzpicture}

  \caption{Líneas de $\vec E$ para $K=4$, calculadas a partir de $\phi_1$ y
  $\phi_2$. Adentro el campo vale $E_0/2$ y las líneas van con el doble de
  separación que lejos. Las líneas exteriores que llegan a la esfera terminan en
  la carga de polarización negativa del lado izquierdo y vuelven a nacer en la
  positiva del derecho; ninguna llega perpendicular a la superficie.}
\end{figure}

\begin{notebox}[Más separadas, no más apretadas]
En el pizarrón el docente dibuja las líneas interiores «más apretadas». Con $K>1$
el campo interior es \textbf{menor} que $E_0$, así que las líneas de $\vec E$
adentro van \textbf{más separadas}; las que se concentran en la esfera son las de
$\vec D$, porque $D_2 = 3K\varepsilon_0 E_0/(K+2) > \varepsilon_0 E_0$. La figura
dibuja $\vec E$.
\end{notebox}

\begin{warnbox}[El campo no es radial en el borde]
Contra lo que pasaba con la esfera conductora, acá llega «chanfleado», con
componente radial y tangencial. Llegar perpendicular a la superficie es propiedad
de los \textbf{conductores}, y esto es un dieléctrico.
\end{warnbox}

\subsection{La versión astuta: unicidad}
\label{sec:astuta}

«Cuando terminemos van a decir: ah, pero era más fácil.» Se metió todo, se aplicó
«la picadora de carne» y la solución salió. Pero quien ya hizo varios prácticos
sabe que $P_2$, $P_3$, \ldots\ van a dar cero, y puede poner de entrada sólo la
constante y el término de $P_1$ en cada zona.

\begin{keybox}[Por qué vale]
Por el \textbf{teorema de unicidad}. Si con ese ansatz reducido se encuentra una
solución que cumple todas las condiciones de borde, es \textbf{la} solución. Y si
se puso de menos, se nota: \textbf{no se encuentra solución}, y hay que volver
atrás y agregar términos. Es un arte, no una ciencia; a prueba y error, pero
ahorra cuentas.
\end{keybox}

\begin{notebox}[Cuándo no alcanza el dipolo]
Si en lugar de una esfera se pone un elipsoide con los ejes alineados para
conservar la simetría azimutal, el problema es menos simétrico y aparecen otras
componentes multipolares además de la dipolar.
\end{notebox}

\section{Energía electrostática}

En los últimos minutos se arranca el \textbf{capítulo 4}. Desde Física 1 se sabe
que los argumentos energéticos simplifican problemas, y eso es particularmente
cierto para la interacción electrostática, que es \textbf{conservativa}.

\subsection{El trabajo no depende del camino}

El trabajo de la fuerza electrostática sobre una carga $q$ que va de $A$ a $B$
por una curva $C$ es
\[
W_{A\to B}^{(C)} = \int_C \vec F \cdot d\vec l
= q\int_C \vec E\cdot d\vec l
= -q\int_C \grad\phi\cdot d\vec l
= -q\left[\phi(B)-\phi(A)\right].
\]

\begin{keybox}[Trabajo de la fuerza electrostática]
\[
\boxed{\ W_{A\to B} = -q\,\left[\phi(B) - \phi(A)\right]\ }
\]
\textbf{No depende de la trayectoria}: sólo de los potenciales inicial y final.
\end{keybox}

\begin{figure}[H]
  \centering
  % Clase 11 — trabajo de la fuerza electrostatica entre A y B.
% Dos curvas y no una: el contenido de la figura es que el resultado NO depende
% de cual se tome. El elemento dl va tangente a C.
\begin{tikzpicture}[scale=1]

  \coordinate (A) at (0,0);
  \coordinate (B) at (5.2,0.9);

  % curva C
  \draw[figline, line width=1pt] (A) .. controls (1.2,2.4) and (3.6,2.6) .. (B)
    coordinate[pos=0.5] (M);
  \node[figlbl, figblue] at (0.35,1.35) {$C$};

  % curva C' (alternativa)
  \draw[figline, figgray, dashed, line width=0.8pt] (A) .. controls (1.4,-1.6) and (3.8,-1.2) .. (B);
  \node[figlbl, figgray] at (2.8,-1.45) {$C'$};

  % elemento de desplazamiento, tangente a C en M
  \draw[figvecres] (M) -- ++(0.95,0.14);
  \node[figlbl, figred] at ($(M)+(0.5,0.45)$) {$d\vec l$};

  % carga q sobre C
  \node[figpos] at (M) {};
  \node[figlbl] at ($(M)+(0,-0.42)$) {$q$};

  % extremos
  \fill (A) circle (1.8pt) node[left, figlbl] {$A$};
  \fill (B) circle (1.8pt) node[right, figlbl] {$B$};

\end{tikzpicture}

  \caption{Por $C$ o por $C'$, el trabajo que hace la fuerza electrostática sobre
  $q$ entre $A$ y $B$ es el mismo: sólo depende de $\phi(A)$ y $\phi(B)$.}
\end{figure}

\subsection{La energía es del sistema, no de una partícula}

Como no depende del camino, se puede definir una energía electrostática. Pero
---y esto el docente lo insiste también en Física 3--- \textbf{la energía
potencial no está asociada a una partícula}, sino al \textbf{conjunto} de
partículas presentes. Lo que se calcula es cuánto cuesta \textbf{constituir} el
sistema.

Se define la energía potencial del sistema como el \textbf{trabajo que tiene que
hacer un agente externo} para traer las partículas desde el infinito hasta sus
posiciones, con nivel de referencia $U = 0$ cuando todas están infinitamente
alejadas unas de otras.

\begin{notebox}[El signo]
El trabajo de la fuerza electrostática y el del agente externo tienen signos
opuestos. Es lo mismo que con la gravedad: la fuerza apunta hacia abajo, pero la
energía potencial \textbf{crece} al subir. Es la energía que el agente tiene que
aportar para llevar el cuerpo arriba, y la que la gravedad libera si se lo
suelta.
\end{notebox}

\subsection{Energía de \texorpdfstring{$N$}{N} cargas puntuales}

Sección 4.1: se arma el sistema trayendo las cargas \textbf{de a una}.
\begin{itemize}
  \item \textbf{Carga 1:} trabajo nulo. No hay ninguna otra en la vuelta, así que
    no hay fuerza.
  \item \textbf{Carga 2:} interactúa sólo con la 1:
    $\displaystyle W_2 = \frac{1}{4\pi\varepsilon_0}\frac{q_1 q_2}{r_{12}}$.
  \item \textbf{Carga 3:} interactúa con la 1 \textbf{y} con la 2:
    $\displaystyle W_3 = \frac{1}{4\pi\varepsilon_0}
    \left(\frac{q_1 q_3}{r_{13}} + \frac{q_2 q_3}{r_{23}}\right)$.
  \item Y así sucesivamente.
\end{itemize}

\begin{keybox}[Energía de un sistema de cargas puntuales]
\[
\boxed{\ U = \frac{1}{4\pi\varepsilon_0}
\sum_{i=1}^{N} \sum_{j=1}^{i-1} \frac{q_i\, q_j}{r_{ij}}\ }
\]
\end{keybox}

\begin{warnbox}[Parejas, no «una con todas»]
Son \textbf{todas las parejas posibles}, y cada pareja se cuenta \textbf{una sola
vez} ---el par $(1,2)$ no se distingue del $(2,1)$---.
\end{warnbox}

\medskip
\noindent\emph{La clase siguiente retoma esta fórmula, la reescribe en términos de
los potenciales en cada carga y la extiende a distribuciones continuas y a la
densidad de energía en el campo.}

\end{document}
